\section{Localized boundary amplitudes and spatial radial evolution}
\label{sec:localized-response}

The auxiliary product in Section~\ref{sec:localization} prescribes its
modes from arithmetic. We now test an independently specified free
localization problem. It supplies a one-sided gamma amplitude and an
explicit scalar-endpoint state, but no identification with the spatial
transfer. The obstruction occurs before fitting the pole factors:
ordinary spatial dilation fails to preserve the fixed protected
representatives. Canonical Hodge representatives allow a different
radial evolution, whose elementary scalar channel contains only the
lowest spatial mode.

The localization and Hodge inputs are established results of
\cite{DPYHiggs,DedushenkoGluing}. Their application and the comparison
below are internal research calculations, without independent
specialist review. The statements concern the specified free sector,
not every supersymmetric defect or Wilson construction.

\subsection{A specified charge, boundary state and circle control}
\label{subsec:localized-boundary}

Take one free three-dimensional $\mathcal N=4$ hypermultiplet on a
round hemisphere $HS^3$ of radius $r$. Use the Higgs boundary
polarization of \cite[Sections~4.2 and 4.3.2]{DedushenkoGluing}.
In that reference's sphere-algebra notation the charge is
\begin{equation}\label{eq:localized-charge}
 \mathcal Q^H=Q_1^++Q_2^-
 =\mathcal Q_1^{\ell+}+\mathcal Q_1^{r-}
  +\mathcal Q_2^{\ell-}+\mathcal Q_2^{r+}.
\end{equation}
The $\ell,r$ superscripts here label supersymmetry algebras. They do
not denote spatial angular momentum or an additional radius.
For this equivariant charge, the protected complex is understood on
the $\mathcal Q^2$-invariant subspace; unrestricted nilpotence is not
assumed.
The scalar boundary data and auxiliary completion are
\begin{equation}\label{eq:localized-boundary-data}
 q_1|=Y,\qquad \widetilde q_2|=\bar Y,\qquad
 G|=-\frac{\ii}{2r}Y,\qquad
 \bar G|=\frac{\ii}{2r}\bar Y,
\end{equation}
with the fermionic completion in
\cite[equations~(80)--(82)]{DedushenkoGluing}. In the free theory,
$G=\ii\partial_\perp q_2$, so
$\partial_\perp q_2|=-Y/(2r)$; zero Neumann data would define a
different problem. There is no dynamical gauge multiplet in this
control. An embedding into the four-dimensional Wilson problem is
not assumed.

Put $\epsilon=(4\pi r)^{-1}$ and $x=|Y|^2/\epsilon$. The localized
empty-hemisphere state is
\begin{equation}\label{eq:localized-vacuum}
 \Psi_0(Y,\bar Y)=(\pi\epsilon)^{-1/2}e^{-x},\qquad
 \int_{\C}|\Psi_0|^2\,\mathrm d^2Y=\frac12.
\end{equation}
For a neutral radial state
$\Psi=(\pi\epsilon)^{-1/2}f(x)$, the normalized boundary
Fourier--Mellin transform has radial part
\begin{equation}\label{eq:localized-mellin}
 \mathcal Mf(\sigma)=\frac1{\sqrt{2\pi}}
   \int_0^\infty x^{s-1}f(x)\dd x,\qquad
 s=\frac12-\ii\sigma,\quad \sigma\in\R.
\end{equation}
The angular Fourier label is zero for these states. In particular,
\begin{equation}\label{eq:localized-gamma}
 \mathcal M(e^{-x})(\sigma)=\frac{\Gamma(s)}{\sqrt{2\pi}}.
\end{equation}
The integral converges for $\Ree s>0$ and extends meromorphically.
Its one-sided simple poles are therefore a prediction of this
specified boundary amplitude, without an arithmetic mode prescription.
This agrees with the mirror hemisphere amplitude in
\cite[equations~(135)--(140)]{DedushenkoGluing}.

The circle determinant is a different object. The polarized scalar
fields have twist
$u^a(\varphi)=(\cos(\varphi/2),\sin(\varphi/2))$ and antiperiodic
boundary conditions. At fixed real mass/background $\mu$, their
kinetic operator is $\partial_\varphi+\mu$ and has weights
$\mu+\ii(n+1/2)$, $n\in\mathbb Z$
\cite[Section~6.1]{DPYHiggs}. Pairing both frequency directions gives
\begin{equation}\label{eq:localized-circle}
 \frac{Z_{\rm circ}(\mu)}{Z_{\rm circ}(0)}
 =\prod_{n=0}^\infty
   \frac{(n+1/2)^2}{(n+1/2)^2+\mu^2}
 =\frac1{\cosh\pi\mu},\qquad
 Z_{\rm circ}(0)=\frac12.
\end{equation}
Correspondingly,
$|\Gamma(1/2-\ii\sigma)|^2/(2\pi)
 =[2\cosh(\pi\sigma)]^{-1}$.
The gamma amplitude before gluing and the paired determinant
must not be interchanged. In particular, the direct real-mass
shift ratio tends to $e^{2\pi\omega}$ at large mass, whereas
the arithmetic factor \eqref{eq:small-transfer-factor} behaves
as $(2\pi/p)^\omega$ at large positive $p$.

\subsection{The actual two-scalar state and its adjoint}
\label{subsec:localized-pair}

The boundary operators act by
\begin{equation}\label{eq:localized-operators}
 Q_L=Y,\qquad \widetilde Q_L=-\epsilon\partial_Y,\qquad
 \widetilde Q_R=\bar Y,\qquad Q_R=-\epsilon\partial_{\bar Y},
 \qquad [\widetilde Q_L,Q_L]=-\epsilon I.
\end{equation}
The right operators use the opposite ordering of the right action
\cite[equations~(125)--(126)]{DedushenkoGluing}.
For the free massless semicircle, separated opposite-pole scalar
endpoints with identity transport give
\begin{align}
 \Psi_{\rm pair}
 &=Q_L\widetilde Q_R\Psi_0
   =Y\bar Y\Psi_0=\epsilon x\Psi_0,
 \label{eq:localized-pair-state}\\
 \widehat\Psi_{\rm pair}
 &=\frac{\epsilon}{\sqrt{2\pi}}\Gamma(s+1).
 \label{eq:localized-pair-gamma}
\end{align}
The normalized gluing overlap is $\epsilon/2$, agreeing with the
ordered propagator
$-\epsilon\,\operatorname{sgn}(\varphi_1-\varphi_2)/2$.
No contact subtraction is forced at separated endpoints. Reversing
the coincident ordering gives $\epsilon(x-1)\Psi_0$, and choosing
the symmetric local composite gives $\epsilon(x-1/2)\Psi_0$.
These are distinct observables; none forces an additional quartic
insertion into \eqref{eq:localized-pair-state}.

The real-coordinate boundary norm reduces to $L^2((0,\infty),\dd x)$.
For $N=-x\partial_x$, integration by parts and Mellin Plancherel give
\begin{equation}\label{eq:localized-field-adjoint}
 N^*=1-N,\qquad
 \mathcal M(Nf)=s\,\mathcal Mf,\qquad
 \norm{f}_{L^2(\dd x)}^2
 =\int_{\R}|\mathcal Mf(\sigma)|^2\dd\sigma.
\end{equation}
These identities hold first on a suitable smooth core. The centered
operator $N-1/2$ is skew-adjoint and generates unitary dilation of
the \emph{field magnitude}. It has not been identified with spatial
dilation. On the real contour ordinary conjugation sends $s$ to
$1-s$. For example,
\[
 \int_0^\infty |xe^{-x}|^2\dd x=\frac14,\qquad
 \int_{\R}\frac{|s|^2}{2\cosh\pi\sigma}\dd\sigma=\frac14,
 \qquad
 \int_{\R}\frac{s^2}{2\cosh\pi\sigma}\dd\sigma=0.
\]
Reusing a vacuum gluing phase with the same analytic polynomial on
both sides would make the last, incorrect norm substitution.
This boundary-coordinate identity also does not prove an isometry
from the full physical radial state space or arithmetic storage.

\subsection{Why ordinary spatial dilation fails on fixed representatives}
\label{subsec:localized-nondescent}

The hemisphere cut is at great-circle angles $\varphi=0,\pi$.
For concentric radial slices, stereographically project from the
antipode of its center. On the protected diameter,
\begin{equation}\label{eq:localized-centered-frame}
 t=2r\tan\frac{\varphi-\pi/2}{2},\qquad
 \Omega=\left(1+\frac{t^2}{4r^2}\right)^{-1},\qquad
 \frac{u(\varphi)}{\sqrt{\Omega}}
 =\frac1{\sqrt2}\left(1-\frac{t}{2r},1+\frac{t}{2r}\right).
\end{equation}
The hemisphere is the ball of radius $2r$. A constant H-basis
rotation writes the protected scalar as
\begin{equation}\label{eq:localized-protected-family}
 O(t)=q'_1(t)+\frac{t}{2r}q'_2(t),\qquad
 \mathcal Qq'_2(t)=\chi(t),\qquad
 \mathcal Qq'_1(t)=-\frac{t}{2r}\chi(t).
\end{equation}
Here $\chi$ is a nonzero free-fermion combination. These relations
are the affine conformal Killing-spinor variation underlying the
twist in \cite[equations~(3.1)--(3.5)]{DPYHiggs}. The basis change
is passive and does not replace \eqref{eq:localized-boundary-data}.

\begin{proposition}[Failure on fixed protected representatives]
\label{prop:localized-nondescent}
At fixed finite $r$, ordinary spatial dilation $D$ does not preserve
all closed representatives of the specified free Higgs complex.
In particular it does not directly induce the field-space Mellin
evolution \eqref{eq:localized-field-adjoint}.
\end{proposition}
\begin{proof}
The scalar has conformal dimension $1/2$. Dilation acts on the
fields, keeping the coefficients of the specified insertion fixed.
Thus
\begin{equation}\label{eq:localized-nonclosure}
 [D,O(t)]=(t\partial_t+\tfrac12)O(t)
                   -\frac{t}{2r}q'_2(t),\qquad
 \mathcal Q[D,O(t)]=-\frac{t}{2r}\chi(t).
\end{equation}
The derivative differentiates the complete protected family.
The last expression gives a nonzero witness at a nonzero interior
point. Equivalently, finite dilation leaves the residual
$(1-e^a)tq'_2(e^at)/(2r)$, apart from the common factor $e^{a/2}$.
Acting on the conformal vacuum gives the same boundary-state
counterexample. A neutral separated pair has two such fermionic
terms; setting one insertion at $t=0$ leaves the other witness.
Hence the failure is preservation of closed representatives,
not merely a nonzero commutator between $D$ and $\mathcal Q$.
\end{proof}

The compensation $\widehat D=D-R_H$, with adapted H weights
$(1/2,-1/2)$ on the doublet, satisfies
$[\widehat D,O(t)]=t\partial_tO(t)$ and is the exact twisted
dilation on the relevant local cohomology
\cite[equation~(3.4)]{DPYHiggs}. It acts trivially there.
The corresponding free protected two-point function is proportional
to an ordering sign: the polarization determinant cancels
$|t_1-t_2|^{-1}$. It does not retain the stationary untwisted
spatial covariance of Section~\ref{sec:free}.

Transporting the charge itself instead relates different complexes:
in the standard flat frame,
$e^{aD}\mathcal Q_r e^{-aD}=e^{a/2}\mathcal Q_{e^ar}$.
That alternative needs specified intertwiners and ordinary adjoints.
Proposition~\ref{prop:localized-nondescent} does not exclude it,
an enlarged unprotected state space, another supersymmetric sector,
or the canonical representative construction below.

\subsection{Canonical Hodge evolution and the lost spatial tower}
\label{subsec:localized-hodge}

For either nilpotent Higgs charge in the free conformal theory,
the positive Hodge operator is proportional to
$B_H=D-R_H\succeq0$. Its kernel gives the common canonical
representatives \cite[equation~(3.4)]{DPYHiggs}. In the joint
energy/R-charge basis, $D$ commutes with
$P_H=1_{\{0\}}(B_H)$. Consequently
\begin{equation}\label{eq:localized-hodge-flow}
 T_H(u)=e^{-uD}|_{\ker B_H},\qquad u\geq0,
\end{equation}
is a contraction semigroup and $D=R_H$ on this space.
Choosing these representatives first is a different operation from
the failed descent on arbitrary representatives. It does not make
$D$ zero or prove an isometry with the hemisphere coordinate norm.

\begin{proposition}[Elementary spatial channel after Hodge projection]
\label{prop:localized-hodge-loss}
For a highest-H-weight elementary free scalar, Hodge projection
retains only its lowest spatial harmonic. In the parity-even axial
covariance channel, the retained and omitted terms are
\begin{equation}\label{eq:localized-omitted-tower}
 C_H(u)=e^{-u/2},\qquad
 C_{\rm omit}(u)=\frac{e^{-5u/2}}{1-e^{-2u}},\qquad u>0.
\end{equation}
\end{proposition}
\begin{proof}
The conformally rescaled scalar cylinder operator is
$-\partial_u^2-\Delta_{S^2}+1/4$. On a one-particle harmonic,
\[
 D=\ell+\tfrac12,\qquad R_H=\tfrac12,\qquad B_H=\ell.
\]
Only $\ell=0$ lies in the Hodge kernel. The parity-even observable
channel in Section~\ref{sec:free} is
\[
 C_{\rm ev}(u)=\tfrac12[C(u,1)+C(u,-1)]
 =\sum_{n\geq0}e^{-(2n+1/2)u}
 =\frac{e^{-u/2}}{1-e^{-2u}}.
\]
Subtracting its $\ell=0$ term proves
\eqref{eq:localized-omitted-tower}.
\end{proof}
The full spatial multiplicity remains $2\ell+1$; the unit coefficients
above describe the selected channel. The aligned and antipodal
observations are used at $u>0$, not as a normalized vector at $u=0$.
Their omitted fraction is $e^{-2u}$, tending to one as $u\downarrow0$,
and, for $\Ree p>1/2$ in the arithmetic comparison,
\begin{equation}\label{eq:localized-omitted-comparison}
 C_{\rm omit}(u)=\frac1{2u}-\frac34+O(u),\qquad
 2\int_0^\infty ue^{-pu}C_{\rm omit}(u)\dd u
 =2\sum_{n\geq1}\frac1{(p+2n+1/2)^2}.
\end{equation}
This is exactly the shape of the positive tower in
\eqref{eq:mode-count-test}. The full arithmetic derivative still
contains $-2/(p-1/2)^2$. Equation~\eqref{eq:localized-omitted-comparison}
does not count the discarded tower as arithmetic positivity.
In particular, removing the lowest scalar mode after this projection
leaves no elementary spatial tower at all.

\subsection{Composite-state characters and endpoint spectral weights}
\label{subsec:localized-character}

The many-particle protected space contains more states. Let $A,B$
denote independent highest-H-weight scalar creators of flavor
charges $+1,-1$, not ordinary adjoints. The free Higgs chiral ring
is polynomial in these generators
\cite[Section~4.3.2]{DedushenkoGluing}. Normalized monomials
$A^aB^b$ have energy $(a+b)/2$ and flavor $a-b$, so
\begin{equation}\label{eq:localized-character}
 \chi_H(u,y)
 =\sum_{a,b\geq0}e^{-u(a+b)/2}y^{a-b}
 =\frac1{(1-ye^{-u/2})(1-y^{-1}e^{-u/2})}.
\end{equation}
At flavor one, the states $A^{k+1}B^k$ have energies $k+1/2$.
Selecting even $k$ gives the ladder $2n+1/2$; additionally deleting
the ground state gives the character $C_{\rm omit}$. These are
additional selections of composite states, not the physical
single-particle angular projection.

For an elementary normalized endpoint ket $E|0\rangle=|1,0\rangle$,
the actual ordinary-adjoint radial response is instead
\begin{equation}\label{eq:localized-endpoint-response}
 \langle0|E^*T_H(u)E|0\rangle=e^{-u/2},\qquad
 \int_0^\infty e^{-pu}\langle0|E^*T_H(u)E|0\rangle\dd u
 =\frac1{p+1/2}.
\end{equation}
A finite polynomial endpoint gives only finitely many energies
in this free control. An infinite response tower needs a specified
observable or coupling and its matrix elements. A matching counting
character does not supply them. This physical radial adjoint is
also not an assertion that a tilded cohomological endpoint is the
ordinary adjoint, or that \eqref{eq:localized-pair-gamma} has acquired
a spatial spectral interpretation.

\subsection{Scope of the arithmetic comparison}

The complete auxiliary arithmetic amplitude is
\begin{equation}\label{eq:localized-arithmetic-amplitude}
 F_{\rm ar}(p)=\pi^{-p/2}(p-\tfrac12)\Gamma(p/2+5/4),
 \qquad
 K_\omega^<(p)=\frac{F_{\rm ar}(p-\omega)}{F_{\rm ar}(p+\omega)}.
\end{equation}
For the trial identification $s=(p+1/2)/2$, the polynomial
$N(2N-1)$ acting on $e^{-x}$ gives
$(2x^2-3x)e^{-x}$ and Mellin factor
$s(2s-1)\Gamma(s)$. This matches the non-exponential part of
\eqref{eq:localized-arithmetic-amplitude}. It is an inverse fit,
not the insertion selected by the separated scalar pair.
Moreover the different trial offset $s=(p+5/2)/2$ changes the
required polynomial to $2N-3$. A degree obstruction based on
one trial map is therefore conditional, not universal.

The Mellin coordinate in \eqref{eq:localized-mellin} is the magnitude
of a boundary field. The spatial variable is logarithmic radius.
The fixed-representative obstruction and the Hodge channel calculation
provide no identification of these variables, their measures or
their adjoints. Analytic continuation does not supply that missing
intertwining relation, and a rescaling chosen to manufacture
$\pi^{-p/2}$ does not explain the fixed arithmetic finite part.
Nor can the coefficient of a $Q$-exact localizing deformation
become the physical shift.

The useful physical alternatives must retain the unprotected
angular modes, specify an enlarged or interface sector, or derive
an observable coupling to the required composite states. They
must then reproduce the causal response and its ordinary adjoint.
The gamma amplitude established here explains a possible special-function
structure; it supplies neither the complete arithmetic transfer nor
an independent cumulative positive norm.
